\documentclass{article}

\usepackage{amsmath,amssymb}
\usepackage{xurl}
\usepackage[colorlinks=true,linkcolor=blue,citecolor=blue,urlcolor=blue]{hyperref}

% Shared commands for notes in docs/paper-gaps/.

\DeclareUnicodeCharacter{2080}{\ifmmode{}_{0}\else\textsubscript{0}\fi}
\DeclareUnicodeCharacter{2081}{\ifmmode{}_{1}\else\textsubscript{1}\fi}
\DeclareUnicodeCharacter{2082}{\ifmmode{}_{2}\else\textsubscript{2}\fi}
\DeclareUnicodeCharacter{2099}{\ifmmode{}_{n}\else\textsubscript{n}\fi}

\newcommand{\Lean}{\textsc{Lean}}
\ExplSyntaxOn
\NewDocumentCommand{\leanid}{m}{
  \ifmmode
    \textsf{\detokenize{#1}}
  \else
    \group_begin:
    \urlstyle{sf}
    \tl_set:Nn \l_tmpa_tl {#1}
    \tl_replace_all:Nnn \l_tmpa_tl {\_} {_}
    \exp_args:NV \path \l_tmpa_tl
    \group_end:
  \fi
}
\ExplSyntaxOff
\providecommand{\tr}{\operatorname{tr}}
\newcommand{\tnleanIssueBase}{https://github.com/LionSR/TNLean/issues/}
\newcommand{\tnleanPrBase}{https://github.com/LionSR/TNLean/pull/}
\newcommand{\tnleanDocBase}{https://sirui-lu.com/TNLean/docs/}
\newcommand{\ghissue}[1]{\href{\tnleanIssueBase#1}{GitHub issue \##1}}
\newcommand{\ghpr}[1]{\href{\tnleanPrBase#1}{PR \##1}}
\newcommand{\doclink}[1]{\href{\tnleanDocBase#1.html}{\path{#1}}}

% Dirac notation and the identity operator, as in blueprint/src/macros/common.tex.
% \providecommand keeps note-local definitions authoritative.
\usepackage{dsfont}
\providecommand{\ket}[1]{|#1\rangle}
\providecommand{\bra}[1]{\langle #1|}
\providecommand{\braket}[2]{\langle #1 | #2 \rangle}
\providecommand{\ketbra}[2]{|#1\rangle\!\langle #2|}
\providecommand{\Id}{\mathds{1}}
\providecommand{\one}{\mathds{1}}
\providecommand{\C}{\mathbb{C}}
\providecommand{\MN}[1]{M_{#1}(\C)}
\providecommand{\MD}{M_D(\C)}

% External referencing for paper sources.  The build helper writes sanitized
% auxiliary files under build/paper-aux/ and excludes duplicated source labels.
\usepackage{xr}

% Import a generated paper auxiliary file with a caller-supplied label prefix.
% The candidate paths support builds from docs/paper-gaps/, the repository root,
% and build/paper-aux/ itself.
\newcommand{\importpaperaux}[2]{%
  \IfFileExists{../../build/paper-aux/#2.aux}{%
    \externaldocument[#1]{../../build/paper-aux/#2}%
  }{%
    \IfFileExists{build/paper-aux/#2.aux}{%
      \externaldocument[#1]{build/paper-aux/#2}%
    }{%
      \IfFileExists{#2.aux}{%
        \externaldocument[#1]{#2}%
      }{}%
    }%
  }%
}

% General external reference macros
\newcommand{\paperref}[2]{\ref{#1-#2}}
\newcommand{\papereqref}[2]{(\ref{#1-#2})}

% CPSV16 source (1606.00608 / MPDO-22-12-17-2.tex) external document and shortcuts
\importpaperaux{cpsv16-}{1606.00608/MPDO-22-12-17-2-tnlean}
\newcommand{\cpsvref}[1]{\paperref{cpsv16}{#1}}
\newcommand{\cpsveqref}[1]{\papereqref{cpsv16}{#1}}

% Verdict marker: \gapnote{<kind>}{<status>}. Kind is the policy
% classification (clarification, local-correction, scope-restriction,
% unfaithful, false-source, open-gap); status is open, wip (elimination
% underway), resolved, or historical. Machine-read by the site index;
% prints a verdict line.
\newcommand{\gapnote}[2]{\par\noindent\textbf{Verdict:} #1 (#2).\par}


% Fallbacks keep this author document compilable when it is built outside its
% source directory and a different command.tex is found first.
\providecommand{\leanid}[1]{\textsf{\detokenize{#1}}}
\providecommand{\doclink}[1]{%
    \href{https://sirui-lu.com/TNLean/docs/#1.html}{\path{#1}}}
\providecommand{\gapnote}[2]{%
    \par\noindent\textbf{Verdict:} #1 (#2).\par}

\title{Model Note: Comparing a Cited Argument with a Formal Statement}
\author{}
\date{\today}

\begin{document}
\maketitle
\gapnote{unfaithful}{open}

% This file is a model for notes in docs/paper-gaps/.
% When making a real note, replace the abstract toy assertion below by the actual
% mathematical assertion, source citation, issue link, and formal declarations.
% The visible document should remain a coherent mathematical note, not a form.
% If the issue is a scalar correction, a missing term, a counterexample, or a
% different formal-library route, specialize the middle sections accordingly.
% If a Mathlib theorem or construction replaces a step from the paper, explain
% that theorem in ordinary mathematics before naming the formal declaration.
% In a real note, the first paragraph should identify the cited assertion and
% the part of the argument in which it is used. Put issue links, pull requests,
% source files, line ranges, labels, and formal declaration names in footnotes
% when they are only traceability data.

\section{The assertion}

This note concerns a hypothetical proposition of the kind used in
Ref.~\cite{Cirac2021Matrix}. The proposition is used later in the same
argument.

Let $\mathcal H$ denote the hypotheses available at the relevant point of the
argument, and let $C$ denote the conclusion needed there. Suppose the cited
source claims
\begin{align}
    \mathcal H &\Longrightarrow C.
    \label{eq:model_source_claim}
\end{align}
In (\ref{eq:model_source_claim}), $\mathcal H$ is the available hypothesis set
and $C$ is the conclusion required by the later argument. A completed note must
name the source theorem, lemma, or equation and cite the source with an existing
bibliography key. Identify a source-paper equation as an equation of the cited
source; do not represent it as a local \verb|\ref|.

\section{The point at issue}

Assume that the source proof invokes a lemma which, in the special case being
used, proves
\begin{align}
    \mathcal H\wedge\mathcal K &\Longrightarrow C,
    \label{eq:model_extra_hypothesis}
\end{align}
where $\mathcal K$ is an additional hypothesis. If the surrounding argument
does not prove $\mathcal K$, then
(\ref{eq:model_extra_hypothesis}) does not justify
(\ref{eq:model_source_claim}). The mathematical question is whether
$\mathcal K$ follows from $\mathcal H$, must be added to the theorem, or can be
avoided by a different argument.

A completed note must show the intermediate step at which $\mathcal K$ enters.
For example, if the argument factors through a statement $D$, write
\begin{align}
    \mathcal H
        &\Longrightarrow D,
    \label{eq:model_first_step}\\
    D\wedge\mathcal K
        &\Longrightarrow C,
    \label{eq:model_missing_step}\\
    \mathcal H\wedge\mathcal K
        &\Longrightarrow C.
    \label{eq:model_derived_claim}
\end{align}
The gap lies between (\ref{eq:model_first_step}) and
(\ref{eq:model_missing_step}), because the first implication does not supply
$\mathcal K$. This calculation locates the missing hypothesis more precisely
than a restatement of (\ref{eq:model_extra_hypothesis}) alone.

\section{The corrected statement}

In this model, suppose that $\mathcal K$ is not known to follow from
$\mathcal H$, but the later proof uses the theorem only when $\mathcal K$ is
available. The corrected statement is
\begin{align}
    \mathcal H\wedge\mathcal K &\Longrightarrow C.
    \label{eq:model_corrected_statement}
\end{align}
The assertion in (\ref{eq:model_corrected_statement}) strengthens the
hypotheses; it does not derive $\mathcal K$ from the original assumptions.

If the formal development proves (\ref{eq:model_corrected_statement}), state
that fact in prose and identify the exact formal declaration in a
footnote.\footnote{Model formal declaration:
    \path{modelCorrected} in \doclink{TNLean/Path/To/File}.}
The formal declaration is evidence for the corrected mathematical statement; it
does not replace its mathematical explanation.

The same principle applies if the formal proof uses a theorem from a library
that is absent from the cited source. State that theorem first as a mathematical
assertion, verify its hypotheses in the notation of the note, and then identify
the corresponding formal declaration in a footnote.

Every substantive display in the completed note must use \verb|align| and carry
a stable semantic \verb|\label{eq:...}| identifier. Refer to local displayed
equations only through \verb|\ref|. Do not use ``the equation above'', raw
local equation numbers, \verb|$$|\ldots\verb|$$|,
\verb|\[|\ldots\verb|\]|, \verb|equation|, or \verb|align*|. Use multiple
aligned lines when an intermediate equality or implication is needed to locate
the gap.

\section{Conclusion}

The verdict in this model is that the original proof uses the additional
hypothesis $\mathcal K$. The corrected statement
(\ref{eq:model_corrected_statement}) suffices for the surrounding argument only
when $\mathcal K$ is supplied. Without a derivation of $\mathcal K$ or a
different proof of $C$, the source claim
(\ref{eq:model_source_claim}) remains unjustified by this route.

\bibliographystyle{alpha}
\bibliography{references}

\end{document}
